\chapter{Rzutowanie typow}

\section{Po co w ogole rzutowac typy}
Rzutowanie jest potrzebne wtedy, gdy dana wartosc ma juz jakis typ, ale w konkretnym
miejscu programu chcesz spojrzec na nia inaczej. Typowe przypadki to:

\begin{itemize}
  \item zamiana napisu \mc{"42"} na liczbe,
  \item zamiana liczby na tekst,
  \item przejscie przez typ \mc{unknown} albo \mc{any},
  \item swiadome dopasowanie danych z zewnatrz do oczekiwanego ksztaltu.
\end{itemize}

MoonChunk wspiera dwa style rzutowania:

\begin{itemize}
  \item styl z nawiasem typu, np. \mc{(int)value},
  \item styl z operatorem \mc{as}, np. \mc{value as int}.
\end{itemize}

\section{Rzutowanie w stylu \mc{(typ)wartosc}}
To zapis zwarty i bezposredni.

\begin{lstlisting}[style=moonchunk,caption={Rzutowanie w stylu C}]
let count: int = (int)"42";
let label: string = (string)15;
\end{lstlisting}

Taki styl jest szczegolnie wygodny, gdy chcesz szybko pokazac, jaki ma byc efekt
konwersji.

\section{Rzutowanie przez \mc{as}}
Druga mozliwosc to zapis z \mc{as}:

\begin{lstlisting}[style=moonchunk,caption={Rzutowanie przez as}]
let count: int = "42" as int;
let text: string = 15 as string;
\end{lstlisting}

Ten styl jest bardzo czytelny w dluzszych wyrazeniach, bo przypomina naturalne czytanie:
\emph{weź wartosc i potraktuj ja jako dany typ}.

\section{Rzutowanie liczb i tekstu}
Najczestsze i najbardziej intuicyjne konwersje to:

\begin{itemize}
  \item tekst do liczby,
  \item liczba do tekstu,
  \item liczba jednego typu do liczby innego typu,
  \item bool do liczby,
  \item liczba do ogolnego typu \mc{number}.
\end{itemize}

\begin{lstlisting}[style=moonchunk,caption={Czeste rzutowania}]
let a: int = "42" as int;
let b: double = "3.5" as number;
let c: string = b as string;
let d: float = (float)10;
\end{lstlisting}

\section{Rzutowanie do \mc{bool}}
To najbardziej specyficzny przypadek w MoonChunk i warto go dobrze zrozumiec.
Bezposredni zapis z \mc{as} do \mc{bool} jest ograniczony.

\begin{lstlisting}[style=moonchunk,caption={Niedozwolone bezposrednie rzutowanie przez as}]
let flag: bool = 3 as bool;
\end{lstlisting}

Taki zapis nie powinien byc akceptowany. Jesli chcesz doprowadzic wartosc do \mc{bool},
masz dwie glówne drogi.

\subsection*{Droga pierwsza: rzutowanie w stylu \mc{(bool)wartosc}}
\begin{lstlisting}[style=moonchunk,caption={Rzutowanie do bool przez zapis nawiasowy}]
let flag: bool = (bool)3;
\end{lstlisting}

\subsection*{Droga druga: przejscie przez \mc{unknown} albo \mc{any}}
\begin{lstlisting}[style=moonchunk,caption={Most przez unknown}]
let flag: bool = 3 as unknown as bool;
\end{lstlisting}

Ten zapis mowi w praktyce: \emph{na chwile potraktuj wartosc jako niejawna, a potem
rzutuj ja do bool}.

\section{Jak zachowuje sie liczba rzutowana do \mc{bool}}
W biezacej semantyce MoonChunk dodatnie wartosci liczbowe po przejsciu do \mc{bool}
są interpretowane jako \mc{true}, a wartosci niedodatnie jako \mc{false}.
Dlatego:

\begin{lstlisting}[style=moonchunk,caption={Przykladowe wyniki rzutowania liczby do bool}]
let a: bool = (bool)3;
let b: bool = (bool)0;
let c: bool = (bool)(-3);
\end{lstlisting}

odpowiednio daje intuicyjnie:

\begin{itemize}
  \item \mc{a = true},
  \item \mc{b = false},
  \item \mc{c = false}.
\end{itemize}

To jedna z tych zasad, ktore warto po prostu zapamietac i stosowac konsekwentnie.

\section{Przechodzenie przez \mc{unknown} i \mc{any}}
W wielu nowoczesnych systemach typow stosuje sie etap posredni, gdy bezposrednia konwersja
jest za mocna albo zbyt ryzykowna. MoonChunk dopuszcza taki most.

\begin{lstlisting}[style=moonchunk,caption={Lancuch rzutowan}]
let visible: bool = (+3) as unknown as bool;
let hidden: bool = (-3) as unknown as bool;
\end{lstlisting}

To nie jest zapis dla poczatkujacego na co dzien, ale warto wiedziec, co oznacza,
poniewaz pojawia sie w bardziej zaawansowanych przypadkach.

\section{Rzutowanie obiektow, slownikow i tablic}
MoonChunk pozwala tez sprawdzac i dopasowywac typy zlozone.

\begin{lstlisting}[style=moonchunk,caption={Rzutowania typow zlozonych}]
let list: array = [1, 2, 3] as array;
let user: object = { name: "Ala" } as object;
let point: dict = { x: 1, y: 2 } as dict;
\end{lstlisting}

Jesli jednak sprobowalibysmy potraktowac liczbe jak tablice albo prostą wartosc jak
obiekt, MoonChunk powinien zglosic blad.

\section{Rzutowanie do \mc{null} i \mc{undefined}}
Sa to przypadki szczegolne. Taka konwersja ma sens tylko wtedy, gdy wartosc faktycznie
jest juz \mc{null} albo \mc{undefined}. MoonChunk nie traktuje ich jako magicznych typow,
ktore mozna dowolnie uzyskac z kazdej wartosci.

\section{Co nie jest rzutowaniem}
Poczatkujacy czasem mieszaja dwa rozne pojecia:

\begin{itemize}
  \item \important{rzutowanie} -- swiadoma zmiana sposobu interpretacji wartosci,
  \item \important{zwykle przypisanie} -- podstawienie wartosci do zmiennej.
\end{itemize}

To, ze zmienna ma typ \mc{bool}, nie znaczy, ze mozna do niej przypisac dowolna liczbe
bez konwersji.

\begin{lstlisting}[style=moonchunk,caption={Przypisanie bez zgodnosci typu}]
let a: bool = 5 > 5;
a = 3;
\end{lstlisting}

To powinno zakonczyc sie bledem typu, bo tu nie zaszlo zadne prawidlowe rzutowanie.

\section{Kiedy rzutowac, a kiedy lepiej nie}
Rzutowanie jest narzedziem bardzo przydatnym, ale naduzywane szybko zamienia kod w cos
trudnego do zrozumienia. Dobra praktyka jest prosta:

\begin{itemize}
  \item rzutuj wtedy, gdy faktycznie przechodzisz miedzy roznymi reprezentacjami danych,
  \item unikaj rzutowan lancuchowych, jesli zwykla deklaracja typu wystarczy,
  \item przy danych zewnetrznych najpierw zrozum, co naprawde przychodzi,
  \item nie traktuj \mc{unknown} i \mc{any} jako domyslnego rozwiazania na wszystko.
\end{itemize}

\section{Rzutowanie w praktyce strony}
Bardzo realistyczny przypadek dotyczy danych z pliku JSON, gdzie niektore wartosci moga
przyjsc jako tekst, a Ty chcesz wykonac na nich obliczenia.

\begin{lstlisting}[style=moonchunk,caption={Konwersja danych zewnetrznych}]
let rawViews: string = "120";
let views: int = rawViews as int;
let text: string = "Wyswietlenia: " + (views as string);
\end{lstlisting}

Dzieki temu jedna wartosc mozna w odpowiednim momencie potraktowac jako liczbe,
a w innym jako tekst gotowy do pokazania na stronie.
